\section{Managerial Process Plan}

The managerial process plan explains how day-to-day and sprint-to-sprint decisions will be made during the project \cite{spp1}. In academic software development, managerial processes are often overlooked because the team is small. However, even a small team requires a clear method for deciding priorities, handling change requests, incorporating feedback, and maintaining communication. Without such a process, valuable time can be lost through confusion, duplication of effort, or unstructured revision.

\subsection{Communication and Coordination}

Team members need regular communication to ensure that development tasks, report writing, and diagram preparation remain aligned \cite{scrum1}. Coordination is especially important in this project because implementation and documentation influence each other. A feature added to the prototype may require changes in the SRS, figures, or planning discussion. Therefore, communication must be continuous and purposeful.

\subsection{Review and Feedback Integration}

Supervisor feedback and internal team review are central managerial mechanisms in this project \cite{agile1}. Feedback may affect the wording of requirements, the focus of the literature discussion, the structure of diagrams, or the priority of certain features. The managerial process must therefore include deliberate review points so that feedback can be integrated systematically rather than casually.

\subsection{Change Management}

Change management refers to the disciplined handling of evolving ideas, corrections, or newly discovered constraints \cite{agile1}. Since the project follows Agile, change is expected. However, not every change should be accepted immediately. The team must judge whether a proposed adjustment improves academic value, strengthens system coherence, and remains feasible within time limits. This process keeps the project flexible without allowing scope expansion to become uncontrolled.

